%%%%%%%%%%%%%%%%%%%%%%%%%%%%%%%%%%
\section{Assumptions and notation}
Let $\Gamma$ the fundamental group of a closed hyperbolic surface. 
Let $\rho, \rho' : \Gamma \rightarrow SL_3\C$ Anosov. 
We use the following notation related to $\rho$ and $\rho'$.
\begin{itemize}
    \item $\xi, \xi' : \partial \Gamma \rightarrow \F^3 $ refers to the continuous, transverse, 
    $\rho$- or $\rho'$- equivariant, dynamics-preserving boundary maps. 
    \item $\Lambda, \Lambda' \subset \F^3$ refers to the limit curves, 
    which are the images of $\xi, \xi'$.
    \item $Th(x) = \alpha_1(x) \cup \alpha_2(x)$ refers to the balanced ideal thickening of $x \in \F^3$.
    \item $Th(\Lambda) = \cup_{x \in \Lambda}Th(x)$ refers to the thickening of the limit curve.     
\end{itemize}

We call $\rho, \rho'$ \emph{topologically conjugate} if there exists a homeomorphism 
$F:\F^3 \rightarrow \F^3$ such that for all $g \in \Gamma$ and all $x \in \F^3$, 
then $F(\rho(g) \cdot x) = \rho'(g) \cdot F(x)$.
We call $F$ the \emph{topological conjugacy} between $\rho$ and $\rho'$.

Recall that $\rho$ is \emph{hyperconvex} if for all $x=(l_x, H_x),y=(l_y,H_y),z=(l_z,H_z) \in \Lambda$, 
then $l_z \not \subset l_x + l_y$ and $(H_x \cap H_y) \not \subset H_z$. 

%%%%%%%%%%%%%%%%%%%%%%%%%%%%%%%%%%%%%%%%%%%%%%%%%%%%%%%%%%%%%%%%%%%%%%%%%%%%%%%%%%
\section{Hyperconvexity is closed under topological conjugacy}

\begin{theorem}\label{thm:hyperconvex-conjugate}
    Hyperconvexity is closed under topological conjugacy. 
    In other words, if $\rho$ and $\rho'$ are topologically conjugate, 
    then $\rho$ is hyperconvex if and only if $\rho'$ is hyperconvex.
\end{theorem}

We use two supporting lemmas to prove this theorem.
First, we show that the relative position of 0-dimensional Richardson varieties of 
$\Lambda$ determines whether $\rho$ is hyperconvex.
\begin{lemma}\label{lem:hyperconvex}
    An Anosov representation $\rho$ is hyperconvex \emph{if and only if} 
    for all $x,y,z \in \Lambda^{(3)}$ we have $R^0_{x,y} \cap R^0_{x,z} = \{x\}$.
\end{lemma}

For the second lemma, assume that
$F:\F^3\rightarrow \F^3$ is a topological conjugacy between $\rho$ and $\rho'$. 
Clearly $F(\Lambda) = \Lambda'$. 
This mapping is pointwise, i.e. for all $\alpha \in \partial \Gamma$, 
then $F(\xi(\alpha)) = \xi'(\alpha)$. 
So for $x \in \Lambda$, we can use the notation $x' := F(x) \in \Lambda'$.

Topological conjugacies may preserve the 0- and 1-dimensional Richardson varieties of 
$\Lambda$ in one of two ways. 
We say $F$ \emph{matches} Richardson varieties if for all $x,y\in \Lambda^{(2)}$, 
then $F(x_y) = x'_{y'}$ and $F(span(x,x_y)) = span(x,x'_{y'})$.
We say $F$ \emph{switches} Richardson varieties if for all $x,y \in \Lambda^{(2)}$, 
then $F(x_y) = x'^{y'}$ and $F(span(x,x_y)) = span(x,x'^{y'})$.
\begin{lemma}\label{lem:richardson}
    For all $x,y \in \Lambda^{(2)}$, 
    then $F (R^0_{x,y}) = R^0_{x',y'}$ 
    and $F (R^1_{x,y}) = R^1_{x',y'}$.
    Furthermore, $F$ either matches or switches Richardson varieties.
\end{lemma}

We also have the following consequence of Lemma \ref{lem:richardson}, 
which is of independent interest:
\begin{corollary}\label{cor:thick}
    For all $x \in \Lambda$, then $F(Th(x)) = Th(x')$.
\end{corollary}

%%%%%%%%%%%%%%%%%%%%%%%%%
\section{Proof of Theorem \ref{thm:hyperconvex-conjugate}.}

We first prove our main theorem using our supporting lemmas.
\begin{proof}(Theorem \ref{thm:hyperconvex-conjugate})
    Suppose $F:\F^3 \rightarrow \F^3$ is a topological conjugacy between 
    Anosov representations $\rho,\rho':\rightarrow SL_3\C$.

    We begin with the forward direction. 
    Assume that $\rho$ is hyperconvex. Let $x,y,z \in \Lambda^{(3)}$.
    By Lemma \ref{lem:hyperconvex}, then $x_y \neq x_z$ and $x^y \neq x^z$.
    By Lemma \ref{lem:richardson}, then the sets
    \begin{align}
        \{F(x_y), F(x^y), F(x_z), F(x^z) \} = \{x'_{y'}, x'^{y'}, x'_{z'}, x'^{z'} \}
    \end{align} contain four distinct points, 
    and in particular $x'_{y'} \neq x'_{z'}$ and $x'^{y'} \neq x'^{z'}$.
    Since $x,y,z$ were general, then this holds for all triples in $x',y',z' \in \Lambda'$, 
    and $\rho'$ is hyperconvex.

    The other direction proceeds similarly. Assume that $\rho$ is \emph{not} hyperconvex. 
    By Lemma \ref{lem:hyperconvex}, there exist $x,y,z \in \Lambda^{(3)}$ such that 
    $x_y = x_z$ or $x^y = x^z$. 
    By Lemma \ref{lem:richardson}, then $F$ either matches or switches Richardson varieties.
    So then $x'_{y'} = x'_{z'}$ or $x'^{y'}=x'^{z'}$. 
    Therefore $\rho'$ is \emph{not} hyperconvex.
\end{proof}
We also prove our lemma linking Richardson 0-dimensional varieties to hyperconvexity.
\begin{proof}(Lemma \ref{lem:hyperconvex})
    Suppose $\rho$ is Anosov.
    Let $x = (l_x, H_x),y = (l_y, H_y),z=(l_z, H_z) \in \Lambda^{(3)}$.
    Since $x,y,z$ are pairwise transverse, then 
    \begin{align}
        y_x, y^x, z_x, z^x \not \in R^0_{x,y} \cap R^0_{x,z}.
    \end{align}
    It suffices to check the pairs $x_y,x_z \in \alpha_1(x)$ and $x^y,x^z \in \alpha_2(x)$.
    
    By definition, $x_y = (l_x, l_x + l_y)$ and $x_z = (l_x, l_x + l_z)$, 
    which are distinct if and only if $l_x + l_y \neq l_x + l_z$, 
    or equivalently if and only if $l_z \not \subset l_x + l_y$.

    Similarly, $x^y = (H_x \cap H_y, H_x)$ and $x^z = (H_x \cap H_z, H_x)$, 
    which are distinct if and only if $H_x \cap H_y \neq H_x \cap H_z$, 
    or equivalently if and only if $H_x \cap H_y \not \subset H_z$.

    So $R^0_{x,y} \cap R^0_{x,z} = \{x\}$ for all $x,y,z \in \Lambda^{(3)}$ 
    is equivalent to $\rho$ hyperconvex. 
\end{proof}


%%%%%%%%%%%%%%%%%%%%%%%%%%%%%%%%%%%%%%
\section{Action of $\rho(g)$ on $\F^3$}
The remainder of this chapter is devoted to the proof of Lemma \ref{lem:richardson}.

In this section we describe the action of $\rho(g)$ on $\F^3$ for $g \in \Gamma$ infinite order. 
Since $\rho$ is Anosov, then $\rho(g)$ is strongly loxodromic, 
i.e. its eigenvalues have distinct magnitude.
That means $\rho(g)$ has attracting and repelling limit curve points 
$x = \xi(g^+)$ and $y = \xi(g^-)$, respectively. 

Under the action $\rho(g)$, most of $\F^3$ is attracted to $x$ and repelled by $y$. 
We define the \emph{uphill} of $x$ as the subset attracted to $x$ under $\rho(g)$:
\begin{align}
    U_{\rho(g)}(x) := \{f \in F^3 : \rho(g^n) \cdot f \rightarrow x \}.
\end{align} 
Similarly, we define the \emph{downhill} of $y$ as the subset repelled by $y$ under $\rho(g)$, 
or equivalently the subset attracted by $y$ under $\rho(g^{-1})$:
\begin{align}
    D_{\rho(g)}(y) = U_{rho(g^{-1})}(y):= \{f \in \F^3 : \rho(g^{-n}) \cdot f \rightarrow y\}.
\end{align}
We then define the \emph{exceptional set} of $\rho(g)$ as the subset 
\emph{not} attracted to $x$ or \emph{not} repelled by $y$ under $\rho(g)$:
\begin{align}
    E_{\rho(g)} = \F^3 \setminus (U_{\rho(g)}(x) \cap D_{\rho(g)}(y)).
\end{align}

The set $E_{\rho(g)}$ is nonempty and closed, since it is complement of an open set. 
It contains the fixed points of $\rho(g) \cdot \F^3$, which are precisely $R^0_{x,y}$.
It suffices to show this for $x,y,\rho(g)$ as in Example \ref{eg:r0}, 
since up to conjugation $\rho(g) = diag(\lambda, \mu, \nu)$.

We can partition $E_{\rho(g)}$ using each element's attracting and repelling fixed points.
By taking the closure of these subsets, we recover the 1-dimensional Richardson varieties $R^1_{x,y}$.
\begin{itemize}
    \item $span(x, x_y) = \alpha_1(x) = \overline{U(x) \cap D(x_y)}$
    \item $span(x, x^y) = \alpha_2(x) = \overline{U(x) \cap D(x^y)}$
    \item $span(x_y,y_x) = \overline{U(x_y) \cap D(y_x)}$
    \item $span(x^y,y^x) = \overline{U(x^y) \cap D(y^x)}$
    \item $span(y_x, y) = \alpha_1(y) = \overline{U(y_x) \cap D(y)}$
    \item $span(y^x, y) = \alpha_2(y) = \overline{U(y^x) \cap D(y)}$
\end{itemize}
So $E_{\rho(g)}$ is homeomorphic to the union of six copies of $\CP^1$, 
each identified with two neighbors in a cyclic order $(x, x_y, y_x, y, y^x, x^y)$.

% TODO add a figure showing the 0- and 1- dimensional richardson varieties of x,y
%%%%%%%%%%%%%%%%%%%%%%%%%%%%%%%%%%%%%%%%%%%%%%%%%%%%%%%%%%%%%%%%%%%%%%%%%%
\section{$F$ preserves the Richardson 1-varieties of pairs $\xi(g^+), \xi(g^-)$}

Assume $F:\F^3 \rightarrow \F^3$ is a topological conjugacy between 
$\rho,\rho':\Gamma \rightarrow SL_3\C$. 
We use the notation $x' = F(x), y' = F(y)$.

First, we simplify the conclusions that we need to check.
\begin{proposition}\label{prop:simplify}
    Suppose that $F$ a topological conjugacy such that for some $x,y \in \Lambda^{(2)}$
    we have $F$ preserves the Richardson 1-varieties of $x$ and $y$, i.e. $F(R^1_{x,y}) = R^1_{x',y'}$.
    Then
    \begin{itemize}
        \item $F$ preserves the Richardson 0-varieties of $x$ and $y$, i.e. $F(R^0_{x,y}) = R^0_{x',y'}$.
        \item $F$ either matches or switches the Richardson varieties of $x$ and $y$.
    \end{itemize}
\end{proposition}

\begin{proof}
    Let $F$ is a homeomorphism that preserves the Richardson 1-varieties of $x$ and $y$.
    First, let $z \in R^0_{x,y}$. Then $z = U \cap V$ for distinct $U,V \in R^1_{x,y}$.
    So then 
    \begin{align}
        F(z) &= F(U \cap V) \\
        &= F(U) \cap F(V) \text{ since $F$ a homeomorphism, which } \\
        &= U' \cap V' \text{ where } U',V' \in R^1_{x',y'}, \text{ so }\\
        &= z' \in R^0_{x',y'}.
    \end{align}
    So $F(R^0_{x,y}) \subset R^0_{x',y'}$. 
    We can prove the other containment by reversing the argument, so $F(R^0_{x,y}) = R^0_{x',y'}$.
    
    Recall that $x' = F(x), y' = F(y) \in \Lambda'$. 
    Since $x \in span(x,x_y) \cap span(x,x^y)$, then $x' \in F(span(x,x_y)) \cap F(span(x,x^y))$,
    so either $F(span(x,x_y)) = span(x',x'_{y'})$ or $F(span(x,x_y)) = span(x,x^y)$.
    In the first case, $F$ matches the Richardson varieties of $x,y$ and $F(x_y) = x'_{y'}$, 
    and in the second case $F$ switches the Richardson varieties of $x,y$ and $F(x_y) = x'^{y'}$.
\end{proof}

Now, we show that $F$ preserves Richardson 1-varieties. 
We first consider the case where $x,y$ are the attracting and repelling flags of an infinite order element $g \in \Gamma$.
\begin{proposition}\label{prop:x_y}
    Let $x,y \in \Lambda$ such that $x=\xi(g^+)$ and $y = \xi(g^-)$. 
    Then $F$ preserves the Richardson 1-varieties of $x$ and $y$.
\end{proposition}

\begin{proof}
    Since $F$ is a topological conjugacy, it must preserve the exceptional set of $\rho(g)$, 
    so $F(E_{\rho(g)}) = E_{\rho'(g)}$.
    Recall that the exceptional set of $\rho(g)$ is the union of the Richardson 1-varieties of $x$ and $y$,
    that is:
    \begin{align}
        E_{\rho(g)} = \cup_{V \in R^1_{x,y}} V.
    \end{align}

    Since $F|_{E_{\rho(g)}}$ is continuous and injective, then it must send 1-varieties to 1-varieties,
    so $F(R^1_{x,y}) = R^1_{x',y'}$.
\end{proof}

%%%%%%%%%%%%%%%%%%%%%%%%%%%%%%%%%%%%%%%%%%%%%%%%%%%%%%%%%
\section{$F$ preserves Richardson 1-varieties of $\Lambda$}

In this section, we extend our argument from pairs 
$\xi(g^+), \xi(g^-)$ to all distinct pairs $\Lambda^{(2)}$.
The main idea is that pairs of the form $\xi(g^+),\xi(g^-)$ are dense in $\Lambda^2$.

First, we extend Proposition \ref{prop:x_y} to Richardson varieties of pairs $x, z$ where 
$x = \xi(g^+)$ and $z$ is a (distinct) generic point on $\Lambda$.
\begin{proposition}\label{prop:x_z}
    Let $g \in \Gamma$ infinite order, and define $x := \xi(g^+)$. 
    Let $z \in \Lambda \setminus \{x\}$. 
    Then $F$ preserves the Richardson 1-varieties of $x$ and $z$.
\end{proposition}
We observe that Richardson 1-varieties vary continuously on transverse pairs of flags:
\begin{proposition}\label{prop:continuous}
    Let $\{x_i\}_{i=1}, \{y_i\}_{i=1} \in (\F^3)^{\N}$ such that 
    for all $i \in \N$, the pair $(x_i,y_i)$ are transverse.
    Assume $x_i \rightarrow x$ and $y_i \rightarrow y$, 
    and that the pair $(x,y)$ is transverse.
    Then $lim_{i \rightarrow \infty} R^1_{x_i,y_i} = R^1_{x,y}$.
\end{proposition}
This continuity extends to each type of Richardson variety,
e.g. $lim_{i \rightarrow \infty} x_i^{y_i} = x^y$.

We use the following result from group theory:
\begin{remark}[TODO REFERENCE]\label{rmk:qgeo}
    Let $\Gamma$ hyperbolic with boundary $\partial \Gamma \equiv S^1$, 
    and let $x,z \in \partial \Gamma$. 
    Then there exists a quasigeodesic sequence $\{g_i\}_{i=1} \in \Gamma^\N$ 
    such that for all $y \in \partial \Gamma$:  
    \begin{align}
        \\lim_{i \rightarrow \infty} g_i \cdot y = 
        \begin{cases}
            x : y = x\\
            z : y \neq x
        \end{cases}.
    \end{align}
\end{remark}
\begin{proof}(Proposition \ref{prop:x_z})
    If $z = \xi(g^-)$ then we are done by Prop \ref{prop:x_y}, so assume not. 
    Set $y := \xi(g^-)$.

    Let $\{g_i\}_{i=1} \in \Gamma^\N$ as in Rem. \ref{rmk:qgeo}, 
    such that $\lim_{i=1}\rho(g_i)\cdot y = z$ and $\\lim_{i=1}\rho(g_i)\cdot x = x$. 
    For all $i\in \N$, define the following:
    \begin{align}
        x_i := \rho(g_i) \cdot x, \\
        y_i := \rho(g_i) \cdot y, \\
        x'_i := F(x_i), \text{ and } \\
        y'_i := F(y_i).
    \end{align}
    By Prop. \ref{prop:x_y}, then for all $i \in \N$, 
    $F$ preserves the Richardson 1-varieties of $x_i$ and $y_i$.
    By Prop. \ref{prop:continuous}, then $\lim_{i\rightarrow \infty} R^1_{x_i,y_i} = R^1_{x,z}$.
    By continuity of $F$, then $\lim_{i\rightarrow \infty}x'_i = x'$ and $\lim_{i\rightarrow \infty}y'_i = z'$.
    We apply Prop. \ref{prop:continuous} again to the transverse sequence $(x'_i, y'_i)$, 
    so then 
    \begin{align}
        F(R^1_{x,z}) 
        &= \lim_{i\rightarrow \infty} F(R^1_{x_i,y_i}) \\
        &= \lim_{i\rightarrow \infty} R^1_{x'_i,y'_i} \\
        &= R^1_{x',z'}.
    \end{align}
\end{proof}

Next, we extend Proposition \ref{prop:x_z} to Richardson varieties of pairs $w,z$ 
where both $w$ and $z$ are distinct generic points of $\Lambda$.
\begin{proposition}\label{prop:w_z}
    For all $w,z \in \Lambda^{(2)}$, 
    then $F$ preserves the Richardson 1-varieties of $w$ and $z$.
\end{proposition}
\begin{proof}
    If $w$ or $z$ is the attracting fixed point of an infinite order element $g \in \Gamma$, 
    then we are done by Prop \ref{prop:x_z}. Assume not.
    
    Let $\{x_i\}_{i=1} \in \Lambda^\N$ such that each $x_i = \xi(g_i)$ is the 
    attracting fixed point of an infinite order word, and $\lim_{i\rightarrow \infty} x_i = w$.
    This sequence exists because attracting fixed points are dense in $\Lambda$.
    Define $x'_i = F(x_i)$ for all $i\in \N$. 
    By continuity of $F$, then $\lim_{i\rightarrow \infty}x'_i = w'$.

    The proof from here proceeds similarly to that of Prop. \ref{prop:x_z}.
    By Prop. \ref{prop:x_z}, then for all $i \in \N$, 
    $F$ preserves the Richardson 1-varieties of $x_i$ and $z$.
    By Prop. \ref{prop:continuous}, then $\lim_{i\rightarrow \infty} R^1_{x_i,z} = R^1_{w,z}$.
    We apply Prop. \ref{prop:continuous} again to the transverse sequence $(x'_i, z')$, 
    so then 
    \begin{align}
        F(R^1_{w,z}) 
        &= \lim_{i\rightarrow \infty} F(R^1_{x_i,z}) \\
        &= \lim_{i\rightarrow \infty} R^1_{x'_i,z'} \\
        &= R^1_{w',z'}.
    \end{align}
\end{proof}

%%%%%%%%%%%%%%%%%%%%%%%%%%%%%%%%%%%%%%%%%%%%%%%%%%%%%%
\section{$F$ matches or switches Richardson varieties consistently}

In this section, we show that $F$ either matches all Richardson varieties of $\Lambda$ 
or switches all Richardson varieties of $\Lambda$.
We also highlight how $F$ preserves the balanced ideal thickening of each $x \in \Lambda$.

Recall that $Th(x) = \alpha_1(x) \cup \alpha_2(x)$, 
and equivalently $Th(x) = span(x,x_y) \cup span(x,x^y)$ 
for some $y \in \Lambda$ distinct from $x$. 
\begin{proof}(Corollary \ref{cor:thick})
    Let $x \in \Lambda$, and let $y\in\Lambda$ distinct from $x$. 
    By Prop. \ref{prop:w_z}, $F$ preserves the Richardson 1-varieties of $x$ and $y$.
    By Prop. \ref{prop:simplify}, $F$ either matches or switches the Richardson varieties of $x$ and $y$.
    So then 
    \begin{align}
        F(Th(x)) &= F(span(x,x_y) \cup span(x,x^y)) \\
        &= span(x',x'_{y'}) \cup span(x',x'^{y'}) \\
        &= Th(x').
    \end{align}
\end{proof}

\begin{corollary}
    If $F$ matches (resp. switches) the Richardson varieties of some pair of $x,y \in \Lambda^{(2)}$,
    then $F$ matches (resp. switches) the Richardson varieties of all such pairs.
\end{corollary}
\begin{proof}
    Observe that 
    \begin{align}
        \alpha_1(\Lambda) 
        &= \cup_{x \in \Lambda} \alpha_1(x) \\
        &= \cup_{x,y \in \Lambda^{(2)}} span(x,x_y)
        &\cong S^1 \times \CP^1 \subset \F^3.
    \end{align}
    Similarly, $\alpha_2(\Lambda) \cong S^1 \times \CP^1$.
    In particular, $\alpha_1(\Lambda)$ and $\alpha_2(\Lambda)$ are connected subsets of $\F^3$
    whose intersection is $\Lambda$.

    By Cor. \ref{cor:thick}, clearly 
    \begin{align}
        F(\alpha_1(\Lambda)) \cup F(\alpha_2(\Lambda)) 
        &= F(\alpha_1(\Lambda) \cup \alpha_2(\Lambda)) \\
        &= F(Th(\Lambda)) \\
        &= Th(\Lambda') \\
        &= \alpha_1(\Lambda') \cup \alpha_2(\Lambda').
    \end{align}

    Since $F$ is a homeomorphism which sends $\Lambda$ to $\Lambda'$, 
    then either $F(\alpha_1(\Lambda)) = \alpha_1(\Lambda')$,
    or $F(\alpha_1(\Lambda)) = \alpha_2(\Lambda)$.
    
    In other words, if $F$ matches the Richardson varieties of one pair in $\Lambda$, 
    then it matches the Richardson varieties of all such pairs.
    Conversely, if $F$ switches the Richardson varieties of one pair in $\Lambda$,
    then it switches the Richardson varieties of all such pairs.
\end{proof}
This completes the proof of Lemma \ref{lem:richardson}. 

%%%%%%%%%%%%%%%%%%%%%%%%%%%%%%
\section{Conjectures and future steps}
TODO: This should probably go in a different chapter eventually, maybe introduction or conclusions idk.

\begin{remark}
    For distinct $x,y \in \F^3 \setminus \Lambda$, 
    $F$ will not necessarily preserve the 0-dimensional Richardson varieties.
\end{remark}

\begin{remark}
    Theorem \ref{thm:hyperconvex-conjugate} applies to virtual hyperbolic surface groups as well, 
    e.g. $\Gamma = \pi_1(\mathcal{O}^2)$ where $\mathcal{O}^2$ is a hyperbolic 2-orbifold.
\end{remark}